\section{Data Analysis II: Junior Skills and Senior Traits (RQ2, RQ3)}\label{sec:rq23}

Both questions use the same three-step design: (i) univariate comparison of high vs low outcome
groups (Mann--Whitney U, rank-biserial effect size with 2,000-sample bootstrap CI, Holm
correction); (ii) multivariable models (standardised L2-regularised logistic regression and a
depth-$\leq 3$ decision tree with at least 10 people per leaf), scored by repeated stratified
5-fold cross-validation (100 folds in total); (iii) a permutation test (500 label shuffles) of the
cross-validated AUC.

\textbf{A methodological correction.} The regularisation strength was first tuned on AUC. Because
AUC ignores the size of coefficients, it selected very strong shrinkage, and the odds ratios
collapsed towards 1 (e.g.\ SDS conscientiousness OR 1.26 despite Cohen's $d = 1.85$). Tuning on
log-loss fixed this and also improved cross-validated accuracy (JDS 0.817 $\rightarrow$ 0.864;
SDS 0.901 $\rightarrow$ 0.924). All results below use log-loss tuning.

\subsection{RQ2: which skills are associated with a high salary hike for juniors?}

\begin{table}[H]\centering\small
\caption{RQ2 (JDS, $n=139$): group comparison and multivariable association.}\label{tab:rq2}
\begin{tabular}{@{}l c c c c@{}}
\toprule
Skill area & Mean high / low & Rank-biserial [95\% CI] & Holm $p$ & OR per +1 SD [95\% CI]\\
\midrule
Dashboards \& storytelling & 4.85 / 3.81 & \textbf{0.59} [0.45, 0.72] & $<0.001$ & 2.72 [1.77, 4.73]\\
Maths \& stats & 4.71 / 3.83 & \textbf{0.55} [0.40, 0.70] & $<0.001$ & \textbf{3.01} [1.99, 4.60]\\
Coding & 4.64 / 3.85 & 0.49 [0.31, 0.64] & $<0.001$ & 1.67 [1.10, 2.79]\\
AI \& ML & 4.82 / 4.28 & 0.36 [0.19, 0.52] & $<0.001$ & 2.00 [1.23, 3.23]\\
Big data & 3.94 / 3.75 & 0.12 [$-$0.09, 0.32] & 0.217 & 1.76 [1.18, 2.67]$^\dagger$\\
\bottomrule
\end{tabular}
\end{table}

\begin{figure}[H]\centering
\begin{subfigure}{0.49\linewidth}\includegraphics[width=\linewidth]{fig_jds_effects.png}\caption{Effect sizes}\end{subfigure}
\begin{subfigure}{0.49\linewidth}\includegraphics[width=\linewidth]{fig_jds_models.png}\caption{Models vs baseline}\end{subfigure}
\caption{RQ2: effect sizes with bootstrap CIs (left) and cross-validated model performance (right).}\label{fig:rq2}
\end{figure}

\textbf{Models.} The L2 logistic regression reaches accuracy 0.864, balanced accuracy 0.861 and
AUC \textbf{0.904} (fold range 0.77--0.99) against a 0.525 / 0.50 baseline; the tree reaches AUC
0.809. The permutation test gives AUC 0.903 vs a null mean of 0.494, $p = 0.002$.

\textbf{Decision rule.} The tree is easy to read: juniors with storytelling $>4.15$, maths-stats
$>3.70$ and coding $>4.65$ received a high hike in 50 of 56 cases; those with storytelling
$\leq 4.15$ and maths-stats $\leq 4.60$ received a low hike in 26 of 26 cases.

\textbf{Answer to RQ2.} High-hike juniors score higher on four of the five skill areas. The
strongest signals are \emph{communicating results} (dashboards \& storytelling) and
\emph{quantitative foundations} (maths-stats). Big-data skill alone does not separate the groups.

\subsection{RQ3: which traits are associated with senior success?}

\begin{table}[H]\centering\small
\caption{RQ3 (SDS, $n=161$): group comparison and multivariable association.}\label{tab:rq3}
\begin{tabular}{@{}l c c c c@{}}
\toprule
Trait & Mean high / low & Rank-biserial [95\% CI] & Holm $p$ & OR per +1 SD [95\% CI]\\
\midrule
Openness to experience & 48.5 / 33.3 & \textbf{0.77} [0.65, 0.87] & $<0.001$ & 9.09 [4.65, 24.5]\\
Conscientiousness & 53.7 / 35.7 & \textbf{0.76} [0.62, 0.87] & $<0.001$ & 9.78 [5.59, 21.3]\\
Extraversion & 48.9 / 36.9 & 0.57 [0.42, 0.71] & $<0.001$ & 2.78 [1.28, 6.60]\\
Agreeableness & 47.7 / 41.1 & 0.30 [0.09, 0.50] & 0.002 & 1.96 [0.99, 4.40]\\
Neuroticism & 36.1 / 36.3 & 0.07 [$-$0.12, 0.26] & 0.454 & 2.44 [1.57, 4.69]$^\dagger$\\
\bottomrule
\end{tabular}
\end{table}

\begin{figure}[H]\centering
\begin{subfigure}{0.49\linewidth}\includegraphics[width=\linewidth]{fig_sds_effects.png}\caption{Effect sizes}\end{subfigure}
\begin{subfigure}{0.49\linewidth}\includegraphics[width=\linewidth]{fig_sds_models.png}\caption{Models vs baseline}\end{subfigure}
\caption{RQ3: effect sizes (left) and cross-validated model performance (right).}\label{fig:rq3}
\end{figure}

\textbf{Models.} L2 logistic regression: accuracy 0.924, balanced accuracy 0.922, AUC
\textbf{0.963}; tree: AUC 0.938. Permutation test: AUC 0.949 vs null 0.501, $p = 0.002$.

\textbf{Decision rule.} Every senior with openness $\leq 38.5$ was low-success (52 of 52); openness
$>38.5$, conscientiousness $>36.5$ and agreeableness $>40.5$ was high-success in 79 of 83 cases
(Figure~\ref{fig:fig_sds_tree.png}).

\fig[0.95\linewidth]{fig_sds_tree.png}{RQ3 depth-3 decision tree.}

\textbf{Answer to RQ3.} Openness to experience and conscientiousness are by far the strongest
correlates of senior success, followed by extraversion. This agrees with meta-analytic evidence
that conscientiousness predicts job performance across occupations and that openness predicts
training proficiency \citep{barrick1991}.

\subsection{Suppressor effects and caution}

$^\dagger$Big data (JDS) and neuroticism (SDS) show \emph{no} difference on their own but a positive
coefficient once the other variables are held constant. Each is mildly negatively correlated with the
strong predictors (e.g.\ neuroticism vs extraversion $-0.25$). For neuroticism the logistic coefficient
moves from $-0.01$ ($p = 0.94$) alone to 0.87 with conscientiousness and openness, and 1.37
($p = 0.001$) with all five traits. These are conditional associations in a small sample; we report
them but base no recommendation on them.

The SDS outcome separates unusually cleanly (52 of 52 below an openness threshold), which may
indicate curated sample data; findings should be validated on another cohort before being used for
policy. The very large SDS odds ratios partly reflect this separation, so effect sizes and decision
rules are the more reliable summary.
